\section{The complete leading small spectrum}\label{sec:complete-spectrum}

We keep the setting of Section~\ref{sec:transfer}: a compact curve, and a
simple bundle $E$ with pairwise nonisomorphic Jordan--H\"older factors,
written in the normal form of Lemma~\ref{lem:socle-normal-form} with graph
$\Gamma$ and masses $M_i=Vr_i$. We use the weight grading $w$, the tight subgraph
$\mathcal T$, the slacks $\varepsilon_\alpha$, the exponents
$a_\alpha=1+\varepsilon_\alpha$, and the $a$-quotient graphs of
Section~\ref{sec:diagonal-flow}. These depend only on $\Gamma$ and the ranks.
List the small eigenvalues as $\lambda_1\le\cdots\le\lambda_{m-1}$. Assign
exponents to indices so that larger exponents take smaller indices: the
first $d_{a_{\max}}$ indices receive $a_{\max}$, and so on. Write $a_j$ for
the exponent of index $j$.

\begin{theorem}\label{thm:complete-spectrum}
Assume $(\Gamma,M)$ is nondegenerate. For every smooth initial metric on $E$:
\begin{enumerate}
\item $\tau^{a_j}\lambda_j(h(\tau))\to c_j(h)>0$ for every $j<m$.
\item The coefficients with $a_j=1$ are the positive eigenvalues of the
 tight subgraph with conductances $u^*$ and masses $M_i$, with multiplicity.
 They do not depend on the initial metric. The value $\frac12$ occurs with
 multiplicity at least the number of tight components with at least two
 vertices.
\item For every edge the limit $k^\infty_\alpha(h)=\lim\tau^{a_\alpha}
 k_\alpha(g_\tau)$ exists along the diagonal gauges of
 Lemma~\ref{lem:socle-matching}. For $a_j=a>1$, the coefficients $c_j(h)$
 are the positive eigenvalues of the $a$-quotient graph with conductances
 $k^\infty(h)$, with multiplicity. For every $a>1$ with $d_a\ge1$, their sum,
 and hence the largest of them, is unbounded as $h(0)$ ranges over metrics
 on the single bundle $E$.
\item In smooth unitary flow gauges, for each exponent $a$ and each
 distinct value $\rho$ among the coefficients with $a_j=a$, the spectral
 projection onto $\{\lambda_j:a_j=a,\ c_j(h)=\rho\}$ converges in every
 $L^2\to C^k$ norm to the $\rho$-eigenspace of the $a$-quotient graph. This
 eigenspace is embedded in $\mathcal K$ as block-scalar endomorphisms
 constant on the contracted vertices. For $\rho=\frac12$ at $a=1$ it
 contains $\sum_{i\in\mathcal C}w_i\id_{F_i}$ for each tight component $\mathcal C$ with at
 least two vertices.
\end{enumerate}
\end{theorem}

\begin{proof}
\emph{Eigenvalues.} By Theorems~\ref{thm:diag-convergence}
and~\ref{thm:diag-spectrum}, every solution $g$ of the diagonal flow has
$\tau^{a_j}\nu_j(g(\tau))\to\rho_j(g)>0$. The exponent assigned to each
index is the same for all solutions, and $f=\tau^{-a}$ satisfies the shift
condition. Corollary~\ref{cor:coefficient-transfer} gives (1), with
$|\log(c_j(h)/\rho_j(g^{[T]}))|\le\delta_T$. At $a=1$ the value
$\rho_j(g)$ is the same for all $g$, which proves (2).

\emph{Edge limits.} By \eqref{eq:edge-asymptotics},
$\tau^{a_\alpha}k_\alpha(g^{[T]}(\tau))\to k^\infty_\alpha(g^{[T]})$.
Lemma~\ref{lem:instantaneous-edges} puts the lower and upper limits of
$\tau^{a_\alpha}k_\alpha(g_\tau)$ within $e^{\pm2b_T}$ of
$k^\infty_\alpha(g^{[T]})$ for every large $T$. Since $b_T\to0$, the limit
$k^\infty_\alpha(h)$ exists and equals $\lim_Tk^\infty_\alpha(g^{[T]})$. The
positive eigenvalues of each $a$-quotient graph depend continuously on the
conductances, so $c_j(h)$ is the corresponding eigenvalue for $k^\infty(h)$.

\emph{Unboundedness.} Fix $a>1$ with $d_a\ge1$ and constants $\theta_{\mathcal C}$ with
$\sum_{\mathcal C}M_{\mathcal C}\theta_{\mathcal C}=0$. For a large $s_0$, let $g_*$ be the diagonal map with
logarithms $w s_0+\zeta^*+\sum_{\mathcal C}\theta_{\mathcal C}\mathbf 1_{\mathcal C}$, and set
$h(0)=g_*^*h_0$. This is a metric on $E$.
\begin{itemize}
\item The diagonal flow from $g_*$ at time zero is a time shift of the
 solution considered at the end of the proof of
 Theorem~\ref{thm:diag-spectrum}. Hence its constants satisfy
 $Z_{\mathcal C}=\theta_{\mathcal C}+O(e^{4\max|\theta_{\mathcal C}|-\varepsilon_*s_0})$.
\item Its total conductance is
 $\mathfrak e_0\le Ce^{-s_0}+Ce^{4\max|\theta_{\mathcal C}|}e^{-(1+\varepsilon_*)s_0}$.
\item Compare $h$ on $[0,\infty)$ with the corrected diagonal metrics along
 the diagonal flow from $g_*$. The initial discrepancy is $O(\mathfrak e_0)$,
 and Lemma~\ref{lem:layer-energy} bounds the total residual by
 $O(\sqrt{\mathfrak e_0})$. The barriers give
 $e^{-b}\widehat h_{g(\tau)}\le h(\tau)\le e^b\widehat h_{g(\tau)}$ for
 all $\tau\ge0$, with $b=O(\mathfrak e_0+\sqrt{\mathfrak e_0})$ and constants independent of
 $g_*$.
\end{itemize}
Therefore $c_j(h)\in e^{\pm4b}\rho_j(g)$ for every $j$. Take
$s_0\gg4\max|\theta_{\mathcal C}|/\varepsilon_*$. Then
Theorem~\ref{thm:diag-spectrum}(3) shows that the sum of the coefficients
at exponent $a$ takes arbitrarily large values as $(\theta_{\mathcal C})$ varies.

\emph{Spectral subspaces.} Let $u$ be a normalized eigensection with
$a_j=a$ and $c_j(h)=\rho$, and
transport it to $(\db_0+N_{g_\tau},\widetilde h_\tau)$. Write $u=x+y$ with
$x\in\mathcal K$ and $y\perp\mathcal K$ in $h_0$.
\begin{itemize}
\item Changing $\widetilde h_\tau$ to $h_0$ alters pointwise norms by a
 factor $1+o(1)$. So $\lVert D_{g_\tau}u\rVert_{h_0}^2\le\lambda(1+o(1))$
 with $\lambda=O(\tau^{-a})$. By \eqref{eq:form-lower}, which uses the
 cancellation only in $h_0$,
 \[
  \lVert[N_{g_\tau},x]\rVert^2+\lVert y\rVert^2\le C\lambda.
 \]
\item For edges with $a_\alpha<a$, $k_\alpha(g_\tau)\ge c\tau^{-a_\alpha}$
 by $k^\infty_\alpha(h)>0$. So every limit $x_\infty$ of $x$ is constant on
 the contracted vertices. The identifications with the diagonal gauges consist of a map
 tending to the identity followed by a parallel diagonal unitary map, which
 fixes $\mathcal K$ pointwise. Full-projection convergence
 (Theorem~\ref{thm:transfer}) therefore gives $y\to0$ in every $C^k$
 norm.
\item Test the weak eigenvalue equation against a fixed block-scalar $z$
 constant on the contracted vertices. Then
 $\lVert[N_{g_\tau},z]\rVert^2=O(\tau^{-a})$. Replacing $\widetilde h_\tau$
 by $h_0$ costs $o(\tau^{-a})$. In $h_0$ the complementary term is
 $\langle[N_{g_\tau},y],[N_{g_\tau},z]\rangle$, since $D_0^*[N,z]=0$. It is
 at most
 \[
  \lVert N_{g_\tau}\rVert_{L^\infty}\lVert y\rVert\lVert[N_{g_\tau},z]\rVert
  =O(\tau^{-1/2})O(\tau^{-a/2})O(\tau^{-a/2})=o(\tau^{-a}),
 \]
 using $\mathfrak e_{g_\tau}=O(\tau^{-1})$, which follows from the decay bound of
 Lemma~\ref{lem:layer-energy} together with
 Lemma~\ref{lem:instantaneous-edges}.
\end{itemize}
Multiply by $\tau^a$ and pass to the limit, using the edge limits. Edges
with $a_\alpha<a$ do not contribute, since $z$ is constant across them, and
edges with $a_\alpha>a$ contribute $o(1)$. So $x_\infty$ is a
$\rho$-eigenvector of the $a$-quotient graph. The number of eigenvalues in
this group equals the multiplicity of $\rho$. Limits of orthonormal
eigensections remain orthonormal because $y\to0$. Compactness in
$\mathcal K$ and smooth convergence of $y$ give the stated $L^2\to C^k$
convergence. The last assertion of (4) is
Theorem~\ref{thm:diag-spectrum}(2).
\end{proof}

\begin{corollary}\label{cor:metric-asymptotics}
Under the hypotheses of Theorem~\ref{thm:complete-spectrum}, write
$h(\tau)=A_\tau g_\tau^*\widetilde h_\tau$ with the diagonal gauges $g_\tau$
of Lemma~\ref{lem:socle-matching} and $\widetilde h_\tau\to h_0$. The limits
$Z_{\mathcal C}(h)=\lim_TZ_{\mathcal C}(g^{[T]})$ exist, and, up to a common
scalar,
\[
 \log g_\tau=w\log\tau+y_\infty(h)+o(1),\qquad
 y_\infty(h)=\zeta^*+\sum_{\mathcal C}Z_{\mathcal C}(h)\mathbf 1_{\mathcal C}.
\]
\end{corollary}

\begin{proof}
The edge limits give convergence of
$\log d_i-\log d_j-(w_i-w_j)\log\tau$ along every edge. Since the graph is
connected, all such differences converge.
\end{proof}

Haiden, Katzarkov, Kontsevich, and Pandit \cite[Theorem~1.1]{HKKP2018} and
Lam \cite[Theorem~1.1]{Lam2026} describe the metric growth up to a bounded
term. Corollary~\ref{cor:metric-asymptotics} shows that in the
nondegenerate case this bounded term converges. The limits
$Z_{\mathcal C}(h)$ exist because the edge limits converge and the quotient
of $\Gamma$ by the tight components is connected. So the limit depends on the
initial metric only through the constants $Z_{\mathcal C}(h)$.
